\subsection{Mixed thermal images and their lateral continuations}
\label{sec:mixed-time}

The poles of the numerator of \(B_\nu\) are on the rays
\(\arg\omega=\pm3\pi/4\). In the larger open sector
\(|\arg\omega|<3\pi/4\), the gamma ratios still have uniform
Stirling expansions on closed subsectors. Therefore, for any
\(a>0\), the Laplace kernel has two holomorphic continuations
through \(L=-a\). Denote them by \(\Phi_{\nu,\uparrow}\) and
\(\Phi_{\nu,\downarrow}\), according to the side of the slit
from which they are reached. Rotate the frequency tail to
\(-\theta\) or \(+\theta\), respectively, with
\(\pi/2<\theta<3\pi/4\). The factor \(e^{-L\omega}\) decays on
both rays for \(L\) in a neighbourhood of \(-a\).

To calculate their difference, subtract the two deformed contours
and close on the left. This contour is clockwise. The upper pole
locations and residues of \(B_\nu\) are
\begin{align}
 \omega_{+,m} &=-2(1-\ii)(h+m),\notag\\
 \mathcal R_{+,m}
 &=\frac{2(1-\ii)8^\nu}{\Gamma(-\nu)}
   \frac{(\nu+1)_m}{m!}
   \frac{\Gamma(h+\ii(h+m))}{\Gamma(1-h+\ii(h+m))},
 \qquad m=0,1,\ldots ,
 \label{eq:laplace-pole-residue}
\end{align}
and the lower ones are their complex conjugates. Here
\((q)_m=\Gamma(q+m)/\Gamma(q)\). At the upper pole one has
\(h+s_+=-m\) and \(s_-=\ii(h+m)\); using
\(1/\Gamma(-\nu-m)=(-1)^m(\nu+1)_m/\Gamma(-\nu)\) gives
\eqref{eq:laplace-pole-residue}. The cutoff can be taken below the
first pole since changing it adds an entire function. On the outer arc,
choose radii \(2\sqrt{2}(h+m+1/2)\) to stay away from the poles.
The gamma reflection formula reduces each ratio near its pole ray
to a bounded ratio of sines times a power-bounded gamma ratio.
The exponential damping on the left then makes the outer arc vanish.
It follows that
\begin{equation}
 \mathcal D_\nu(a)
 \equiv \Phi_{\nu,\uparrow}(-a)-\Phi_{\nu,\downarrow}(-a)
 =-4\pi\ii\,{\rm Re}\sum_{m=0}^\infty
    e^{a\omega_{+,m}}\mathcal R_{+,m}.
 \label{eq:laplace-jump}
\end{equation}
This sum converges absolutely for \(a>0\). For positive integer
\(\nu\), \(B_\nu\) is a polynomial and \(\mathcal D_\nu=0\).

Consider the upper image
\begin{equation}
 t_{jk}=\ell\tau+\ii s\tau,\qquad \ell=j-k,\quad s=j+k .
 \label{eq:upper-image-time}
\end{equation}
Continue from the thermal strip along a vertical path slightly to the
right or left of \({\rm Re}\,t=\ell\tau\), and then approach the
image from below. A term of \eqref{eq:laplace-image} with
\(j'-k'=\ell\) and \(s'=j'+k'<s\) has
\(L^+_{j'k'}=-(s-s')\tau\) at the image. It is locally holomorphic on
either of the two continuations just constructed. Later terms in
that column and every negative-frequency term have positive real
Laplace argument. Other columns have nonzero imaginary Laplace
argument. Thus the only term singular at the image is
\(\Phi_\nu(\ii(t-t_{jk}))\). On either continuation within the
sector reached from below, the leading behaviour is
\begin{equation}
 G^>(t)\sim
 \frac{C_\nu\Gamma(2\nu+1)}{2\pi}
 \frac{e^{-\ii\ell\varphi_\nu}}{[\ii(t-t_{jk})]^{2\nu+1}}.
 \label{eq:mixed-coefficient}
\end{equation}
The analytic backgrounds depend on the chosen side. At a point
\(t=\ell\tau+\ii y\) away from the image times, the right-minus-left
difference is
\begin{equation}
 G^>_{\mathrm R}(t)-G^>_{\mathrm L}(t)
 =\frac{C_\nu}{2\pi}e^{-\ii\ell\varphi_\nu}
 \sum_{\substack{s'=|\ell|,\,|\ell|+2,\ldots\\ s'\tau<y}}
    \mathcal D_\nu(y-s'\tau).
 \label{eq:image-jump}
\end{equation}
In particular, the mixed image \((1,1)\) at \(t=\ii\beta\) has
the power in \eqref{eq:mixed-coefficient} even when the cut from the
coincidence point passes through it. The indices \(j,k\) in this
calculation count harmonics in the exact spectrum.
The identity \(G^>(-t-\ii\beta)=G^>(t)\) on the thermal strip
maps these continuations to the corresponding lower images.

\begin{table}[t]
 \centering\small
 \begin{tabular}{ccccc}
 \hline
 \((j,k)\) & \(\nu\) & \(\epsilon=0.16\) & \(\epsilon=0.08\) &
 \(\epsilon=0.04\) \\
 \hline
 (1,1) & 0.75 & 6.042e-04 & 9.206e-05 & 1.521e-05 \\
(2,1) & 0.75 & 6.060e-04 & 9.228e-05 & 1.524e-05 \\
(2,1) & 1.25 & 9.056e-05 & 3.771e-06 & 1.046e-07 \\
(2,2) & 1.5 & 1.060e-04 & 6.355e-06 & 3.795e-07 \\
\hline

 \end{tabular}
 \caption{Maximum relative error over the two lateral continuations
 when extracting the leading coefficient at \(t=t_{jk}-\ii\epsilon\).
 Both numerical continuations use the full spectrum.}
 \label{tab:mixed-check}
\end{table}

Table~\ref{tab:mixed-check} compares the scaled correlator on both
sides with \eqref{eq:mixed-coefficient}. The direct contour integrals
check \eqref{eq:laplace-jump} to absolute error below
\(3\times 10^{-9}\) for \(a=0.8,2.5\) and
\(\nu=0.75,1,1.25,1.5\). At \(\epsilon=0.04\) the difference
between the two scaled correlators, divided by \(C_\nu/(2\pi)\),
agrees with the scaled jump in \eqref{eq:image-jump} to absolute
error below \(5\times10^{-13}\) for the images in the table.
